% RESEARCH_DRAFT_SHA256: 6e2e53899b506b39aa47206a76c33fdbf4dad5351c35f6dd416c8e4d1c3b57f0
% Layout and source-reference wrappers expanded for the web renderer.
\documentclass[11pt,a4paper]{article}
\usepackage[margin=25mm]{geometry}
\usepackage{fontspec}
\setmainfont{lmroman10-regular.otf}[BoldFont=lmroman10-bold.otf,
 ItalicFont=lmroman10-italic.otf,BoldItalicFont=lmroman10-bolditalic.otf]
\usepackage{amsmath,amssymb,mathtools}
\usepackage{unicode-math}
\setmathfont{latinmodern-math.otf}
\AtBeginDocument{\let\setminus\smallsetminus}
\usepackage{xurl,hyperref}
\hypersetup{colorlinks=true,urlcolor=blue}
\setcounter{secnumdepth}{0}
\setlength{\parindent}{0pt}
\setlength{\parskip}{5pt}
\setlength{\emergencystretch}{3em}
\begin{document}
\section*{Research-attempt provenance}
This standalone research attempt retains the complete original definition body below,
including its historical source and preparation notes. Those notes describe the
earlier catalogue preparation; the new attempt and its verification are identified
separately after that body. No original ID, catalogue status, or source statement is changed.

\section*{Point Location in 3D Subdivision}\label{5500013}
{\small UnsolvedMath ID 5500013 · AMR-054-0013 · Computer Science · AMR Open Problem Lists}
\subsection{Definitions and mathematical statement}
A polyhedral subdivision of three-dimensional space is a finite
face-to-face complex whose three-dimensional cells cover a prescribed
polyhedral domain (with an exterior cell if needed). Its faces, edges,
vertices, and incidence relations are given explicitly. Let $n$ denote
its total boundary-description complexity; in a triangulated description
this is the number of constant-size boundary faces up to constant factors.

A point-location query supplies $q\in\mathbb R^3$ and asks which cell
contains it. Boundary points may be handled by a fixed deterministic
tie rule, or the query can report an incident cell. The central question
is whether every such subdivision admits a data structure with
\[
                    O(n)\text{ storage},\qquad
                    O(\log n)\text{ worst-case query time}.
\]
Preprocessing must construct the structure from the subdivision; its
cost should be stated separately. The bounds refer to a standard
real-arithmetic/pointer model with constant-size geometric records.

Counting each polygonal face as one unit regardless of an unbounded
number of vertices would conceal input complexity, so that convention
is not used. The requested simultaneous space and query bounds are
stronger than obtaining either bound alone with an unrestricted cost
for the other.
\subsection{Short English statement}
Can arbitrary three-dimensional polyhedral subdivisions support
logarithmic-time point location using only space proportional to their
input complexity?
\subsection{Sources}
\begin{itemize}
\item[S1] ULAM.AI and the UnsolvedMath Contributors, supplied problems.json, record 5500013 (AMR-054-0013), AMR Open Problem Lists. Catalogue identity and source transcription; CC BY 4.0.
\url{https://huggingface.co/datasets/ulamai/UnsolvedMath}.
\item[S2] Open Problems Project - Computational Geometry, Problem 13. Original source indexed by AMR-054-0013.
\url{https://topp.openproblem.net/p13}.
\end{itemize}

\section{Literature and context check}
TOPP Problem 13 was consulted for the simultaneous space/query target. The
page records historical tradeoffs, not a proof that a particular elementary
hierarchy achieves both desired bounds. The following construction audits that
hierarchy without claiming an impossibility theorem for all data structures.
\section{Research attempt: Testing balanced spatial subdivision as a point-location route}
Consider the unit cube subdivided into $m$ closed horizontal slabs
\[
 C_j=[0,1]^2\times[(j-1)/m,j/m],\qquad 1\le j\le m.
\]
Their full face/edge/vertex description has complexity $\Theta(m)$; neighboring
slabs meet face-to-face, and side faces are subdivided at the same levels. A fixed
tie rule assigns a point on a horizontal face to its upper slab, except the top.
For this family, sorting the $m+1$ heights and binary searching gives $O(m)$
storage and $O(\log(m+1))$ queries, including the tie rule. This is a genuine
restricted solution but elementary, not a claimed advance on arbitrary subdivisions.

The attempted general construction instead recursively bisects a bounding box
by planes perpendicular to the $x$ axis, and stores in each leaf the cells that
intersect it. For depth $h$, there are $2^h$ vertical leaf boxes. Every one of the
$m$ slabs has an interior point in every leaf, so the lists contain exactly
$m2^h$ incidences. If $h=\lceil\log_2m\rceil$, a query reaches a leaf in
$O(\log m)$ steps but that leaf still contains all $m$ possible answers, and
the lists use $\Theta(m^2)$ space. If lists are shared to reduce space, their
logical candidate set still has size $m$; the navigation alone has not answered
the query. This proves a precise failure of the inference that balanced geometric
depth ensures balanced answer complexity.

To express the same issue in recurrence form, let $c(Q)$ be the number of
original cells meeting a node box. A binary cut may have
$c(Q_1)=c(Q_2)=c(Q)$, not $c(Q)/2$. Counting objects only at the root and tree
depth therefore cannot justify linear storage. The slab family saturates this
failure at every level of the displayed tree.

Choosing horizontal cuts repairs this family completely; hence the example
does not defeat adaptive separators. What a successful adaptive proof must
control is the total duplication over all cut cells as well as the number of
query tests. Neither follows from balanced coordinates or a small tree depth.
\section{Verification and adversarial audit}
All cells are axis-aligned boxes, so the example satisfies the face-to-face
input convention without hidden high-degree faces. Storage counts include every
cell-to-leaf incidence. $m=1$ is constant size. Query points on slab boundaries
were assigned explicitly. No universal separator lower bound is inferred.
\paragraph{Audit corrections.} A bad hierarchy was not promoted to a lower bound for arbitrary or adaptive data structures; full boundary-description size is counted.
\paragraph{External computation.} None; symbolic arguments and boundary checks only.
\section{Scope and remaining gaps}
The first gap is an adaptive construction with a proved linear duplication
budget and logarithmic query depth for all polyhedral subdivisions. The favorable
slab algorithm and the unfavorable perpendicular hierarchy do not settle existence
of such a data structure in the requested model.
\paragraph{Integration review.} Codex (GPT-6; exact serving revision not exposed) reviewed the written arguments and expanded layout and source-reference wrappers. Repository metadata and the production web renderer were checked using Python 3.13 and JavaScript with MathJax 3.2.2. This model-assisted review is not independent human proof verification; no PDF compilation or formal proof checking is asserted.

\section{Completion Estimate}
\noindent\textbf{COMPLETION: 0\%}
The calculations below the preserved source statement investigate a concrete route,
but establish no substantiated new progress toward the entire original target.
Known special cases and failed sufficient conditions are not credited as a solution.
\section{References}
The original source references and catalogue attribution are preserved in the
Sources section above. The following URLs were actually consulted for this attempt;
retention of an additional original reference does not claim fresh consultation of it.
\begin{itemize}
\item \url{https://topp.openproblem.net/p13} (consulted 2026-10-11). 
\end{itemize}
\section{Model and contributor disclosure}
OpenAI gpt-6-astra, requested reasoning configuration ultra, generated this research
attempt on 2026-10-11. The exact serving revision is not exposed. Codex GPT-6
performed orchestration and repository-format review; its exact serving revision
is also not exposed. The mathematical route, boundary checks and final adversarial
self-audit followed contribution prompts A then B in
\texttt{docs/contribution\_prompts.md}, after reading \texttt{CONTRIBUTING.md}.
Python 3.14.2 assembled the standalone documents; any mathematical computation is
specified in Verification. No theorem prover was used. No human mathematical
authorship, substantial human proof correction, or independent human proof checking
is asserted. Model-assisted review, source retrieval, and structural checks do not
establish correctness or novelty. Full private reasoning transcripts are unavailable
as reproducibility artifacts; the complete arguments, inputs and outcomes actually
used are stated here or in the companion batch audit.
\section{Final status}
Unresolved. This is a research attempt with an explicit failed route or restricted
calculation, not a full solution, not a current-status certification, and not a claim
that the elementary material is new. The catalogue's original source-status label
is preserved independently of this attempt status.
\end{document}
